\chapter{結論與未來研究}
\label{chapter:conclusion}

\section{研究結論}
\label{sec:conclusions}

本研究成功開發並驗證了一套整合物理資訊機器學習、異質圖神經網路、真實都市 GIS 建物擷取與多目標進化演算法的都市設計優化框架，並以 Rhino/Grasshopper 介面實現了從幾何建模到微氣候評估的閉環即時決策系統。以下針對各核心模組之研究結論進行系統性總結：

\noindent\textbf{一、真實都市建物擷取流程之建立}

本研究建立一套三服務融合之 GIS 建物向量化方法，整合 NLSC TOPO01K（樓高標註）、B5000（幾何驗證）與 EMAP（面積驗證）三項政府開放圖資，取代單一仰賴蒙地卡羅隨機採樣之隨機幾何生成方式，使模型得以在真實都市紋理下進行案例驗證，而非僅於程序化生成之封閉解空間內自我評估。

\noindent\textbf{二、異質圖建構之有效性}

本研究提出以五類關係型邊緣（陰影、植被蒸散、對流、語義、毗鄰）建立的異質圖結構，能有效編碼都市行區中建築量體與開放空間之物理因果關係。$N_{\text{obj}}$ 索引偏移機制實現了建物節點與 6,241 個空氣評估節點在單一張量空間中的統一處理，使 RGCN 能以單次前向傳播完成跨節點類型的關係推理。

\noindent\textbf{三、PI-ST-GNN 模型性能}

在 300 個隨機生成都市場景、共計 3,707,154 個預測點次的測試集評估中，PI-ST-GNN 達到：
\begin{itemize}
    \item 決定係數 $R^2 = 0.9957$，遠超代理模型部署門檻 $R^2 > 0.99$；
    \item 均方根誤差 RMSE $= 0.194°\text{C}$（去正規化後之實際 UTCI 單位），遠小於 UTCI 熱壓力等級最小間距（$6°\text{C}$）；
    \item 熱壓力等級分類準確率 98.3\%，確保預測值能可靠地對應至設計決策所需的六等熱壓力分級體系。
\end{itemize}

\noindent\textbf{四、物理資訊損失函數之作用}

消融實驗證實，三項物理資訊損失（輻射一致性 $\mathcal{L}_{\text{rad}}$、時序平滑性 $\mathcal{L}_{\text{temp}}$、風場阻滯 $\mathcal{L}_{\text{wind}}$）在純數值指標上的貢獻並非單調遞增，但其核心價值在於確保預測結果的\textbf{熱力學合理性}，防止模型輸出違反物理常識的「陰影區比日照區更熱」等悖論，提升模型作為建築設計輔助工具的\textbf{可信賴性}。

\noindent\textbf{五、運算效能對即時設計工作流之解鎖效果}

PI-ST-GNN 的推理延遲僅 \textbf{61.5 ms}（RTX 5070 Ti 上測得，200 次平均），相較傳統 CFD 模擬（2--8 小時）效率提升達 $10^5$ 倍量級。此速度量級的突破直接使以下新型設計互動模式成為可能：在 Grasshopper 環境中即時更新全場域 UTCI 熱圖、NSGA-II 100 代演化在 5 分鐘內完成，相較直接呼叫高擬真度模擬引擎之既有 Wallacei 案例常需數週至數月，本研究之最佳化迴圈得以在設計對話的時間尺度內完成。

\noindent\textbf{六、多目標最佳化的設計決策支援與其現行侷限}

以訓練完成之 PI-ST-GNN 作為適應度求值器，NSGA-II 已實際完成一次真實最佳化運行（族群 40、演化 50 代），在全域 UTCI 最小化與綠化率最大化二項目標間求解帕累托前緣；人行動線熱暴露最小化（第三項目標 $f_2$）已於第~\ref{subsec:from_astar_to_walkway}節完成數學定義與推導，取代僅評估單一起訖點路徑的傳統尋路演算法，但其整合進適應度求值器之工程實作列為近期後續工作。真實運行結果同時誠實地揭露：現行二目標函數在缺乏最低開發強度約束下會收斂至近乎零建築量體之退化解（第~\ref{subsec:pareto_morphology}節），此一發現本身即是本研究對「以 GNN 代理模型驅動之多目標最佳化」在實務部署前必須解決之目標函數設計問題的具體貢獻。


\section{研究貢獻與創新點}
\label{sec:contributions}

本研究之學術貢獻與實務創新可歸納為以下四點：

\begin{enumerate}
    \item \textbf{真實圖資交叉驗證之都市建物擷取方法：}
    本研究提出三服務融合（TOPO01K、B5000、EMAP）之建物向量化方法，以粗細線分離、子區塊輪廓萃取與多字元群集光學文字辨識，解決單一圖資來源無法同時提供精確幾何輪廓與樓層屬性之限制，並以跨服務像素重疊率提供每筆建物幾何之信心度評估，使真實都市量體得以系統性地整合進微氣候預測工作流。

    \item \textbf{延伸物理一致之異質時空圖神經網路架構，並補以物理資訊軟約束損失：}
    本研究之 RGCN 關係型聚合公式與 LSTM 暖啟動機制，採用並延伸 \citet{xin2025urbangraph} 之 UrbanGraph 架構設計（詳見第~\ref{subsec:rgcn}、\ref{subsec:lstm}節）。在此基礎上，本研究之貢獻具體體現於三點：其一，將節點集合顯式區分為建物與空氣兩種異質節點類型，並以獨立編碼器分別處理其靜態幾何與動態氣象特徵，區別於原架構之單一節點類型設計；其二，於 UrbanGraph 原有之圖拓撲硬約束（Hard Constraints）之外，另引入三項\textbf{物理資訊軟約束損失}（輻射一致性、時序平滑、風場阻滯），作為訓練過程中的補充監督機制；其三，以\textbf{多步驟時序預測}一次輸出全日 11 時段 UTCI 場，並整合至第~\ref{chapter:results}章之端到端低延遲推理與設計迴圈中。不同於現有研究多以純資料驅動方式預測單一時步的熱舒適指標 \citep{wu2023data}，本研究之上述三項擴展在應用完整性上有所補強。

    \item \textbf{低延遲代理模型實現即時閉環設計：}
    本研究建立了首個推理延遲低至 61.5 ms、並直接嵌入 Rhino/Grasshopper 設計環境的城市微氣候預測系統。此一延遲水準使得 NSGA-II 進化演算法能在合理的設計時間內（$\approx$5 分鐘）完成百代優化，從根本上改變了性能導向設計（Performance-Driven Design）在微氣候評估環節的工作流程，相較既有以 Wallacei 直接呼叫高擬真度模擬引擎之作法，大幅壓縮了最佳化迴圈所需的等待時間。

    \item \textbf{熱舒適人行動線設計之多目標決策框架：}
    本研究將 GNN 預測之 UTCI 場沿設計者定義之人行動線網絡進行時空平均取樣，建立熱舒適人行動線設計評估指標 $\bar{\Phi}_{\text{walkway}}$（式~\eqref{eq:walkway_exposure}），取代僅評估單一起訖點間路徑之傳統尋路演算法，使評估尺度與都市設計決策一致；其作為 NSGA-II 第三項適應度函數之工程整合列為近期後續工作。搭配 5D 違規向量（FAR、BCR、縮退、容納性、碰撞）之法規感知最佳化，本研究以已完整實作之全域 UTCI 與綠化率二目標完成真實最佳化運行，在帕累托前緣上呈現二者間的設計權衡，並誠實揭露現行目標函數在缺乏開發強度下限約束時之退化解風險，為後續研究提供具體、以實證為基礎的改進方向。
\end{enumerate}


\section{研究局限性}
\label{sec:limitations}

儘管本研究在技術實現與應用驗證上取得顯著成果，仍存在以下局限性。各項局限依「限制描述 → 原因 → 改善方向」三段式逐條說明：

\begin{itemize}
    \item \textbf{氣候情境的單一性：}模型對於冬季強烈東北季風、梅雨季連續陰天或颱風侵襲等極端氣候情境之適應能力尚未系統驗證。原因在於訓練資料以台灣夏至（6月21日）單一氣候日為基礎，此一選擇雖能鎖定全年熱壓力最嚴峻之情境以聚焦訓練訊號，卻使模型未接觸過冬季低太陽高度角、高濕陰天等大幅不同之邊界條件組合。改善方向為系統性引入春、夏、秋、冬四季代表性氣候日之 EPW 資料重新訓練或微調（見第~\ref{sec:future_work}節第 3 項）。
    \item \textbf{場地尺度局限：}模型以 $80 \times 80$ 公尺行區為模擬單元，在街廓連棟型配置下可能存在邊界效應誤差。原因在於現有異質圖架構未建立跨行區之連接邊緣，鄰近行區的熱環境交互影響（如建築群峽谷效應引發的背景風速放大）無法透過訊息傳遞機制傳入模型；此一設計取捨是為了控制單一場景之圖規模與訓練成本。改善方向為將異質圖延伸至多行區聯合建模，以行區間界面作為跨行區連接邊緣（見第~\ref{sec:future_work}節第 2 項）。
    \item \textbf{極端熱壓力等級準確率瓶頸：}UTCI $> 46°\text{C}$ 之極端熱壓力等級分類準確率相對較低，此類高熱情境恰為公共健康最需關注之情境。原因為訓練資料中此等級樣本佔比極低（僅出現於正午直射、高 SVF 節點），使模型在此區間之梯度訊號稀疏，屬於典型的類別不平衡問題而非模型結構性缺陷。改善方向為針對高溫情境進行加權採樣或額外資料擴增，提升極端等級樣本於訓練集中之代表性。
    \item \textbf{訓練資料仍以隨機生成場景為主：}第~\ref{sec:building_extraction}節建立之真實建物擷取流程，目前僅用於既有場地案例驗證，尚未回饋為大規模訓練資料集之直接來源。原因在於真實建物擷取流程於本研究階段之產出規模（單一研究場域之案例集）尚不足以支撐訓練所需之場景多樣性，且部分子區塊之樓高辨識仍需人工複核（見第~\ref{subsec:limitations}節第 5 項），限制了自動化規模化的可行性。改善方向為擴大擷取流程之涵蓋範圍並建立真實建物資料庫（見第~\ref{sec:future_work}節第 1 項），使模型面對明顯超出訓練集幾何統計特徵範圍（如超高層建築群、極端窄巷紋理）之真實場地時，預測結果目前仍應視為初步參考。
    \item \textbf{缺乏真實感測器驗證：}本研究之預測精度評估以 EnergyPlus/Ladybug 模擬輸出為基準，尚未以現地 IoT 感測器之實測 UTCI 資料進行交叉驗證。原因在於現有 IoT 感測器網路（第~\ref{subsec:iot_spatial}節）雖提供溫濕度觀測值，用於氣象邊界條件之空間降尺度校準，但並未直接量測 UTCI 所需之平均輻射溫度與風速，無法作為逐點 UTCI 之直接驗證來源。改善方向為部署具備完整四項微氣候變數量測能力之密集感測器網路，收集連續逐時實測資料以進行物理驗證（見第~\ref{sec:future_work}節第 6 項）。
    \item \textbf{現行多目標函數之退化解風險：}第~\ref{subsec:pareto_morphology}節之真實 NSGA-II 運行結果顯示，現行已實作之二目標函數（全域 UTCI 最小化、綠化率最大化）在缺乏開發強度下限約束時，會收斂至 FAR 遠低於訓練分布範圍之近乎零建築量體退化解。原因在於兩項目標之數學形式皆未要求任何建築量體存在——空地本身之預測 UTCI 通常已低，綠化率亦不以建築存在為前提，而現有 5 維違規向量僅檢核上限，未設定最低開發強度。改善方向為引入最低容積率或建築棟數下限之可行性約束，並儘速完成人行動線熱暴露目標 $f_2$ 之工程整合，以三目標之相互制衡降低單一退化方向的可能性（見第~\ref{sec:future_work}節第 7 項）。
\end{itemize}


\section{未來研究建議}
\label{sec:future_work}

基於上述研究成果與局限性，本研究提出以下未來研究方向：

\begin{enumerate}
    \item \textbf{真實建物資料庫回饋訓練集：}
    將第~\ref{sec:building_extraction}節之三服務融合擷取流程擴大應用於新竹市及周邊都市行區，建立真實建物足跡與樓高之大規模資料庫，逐步取代或補充現有隨機生成場景，使模型訓練分佈與真實都市紋理之統計特性更為一致。

    \item \textbf{多行區與街廓尺度延伸：}
    將異質圖架構延伸至多行區聯合建模，以行區間界面（道路境界線、退縮空間）作為跨行區連接邊緣，研究街廓尺度通風廊道、熱島聚集效應等大尺度都市氣候現象。此延伸需要在圖分割（Graph Partitioning）與批次訓練策略上進行適配。

    \item \textbf{季節性與多氣候場景泛化：}
    系統性地引入春、夏、秋、冬四季代表性氣候日之 EPW 訓練資料，以及不同台灣氣候分區（如花東縱谷、高屏熱帶氣候區）之場地條件，建立具備跨時間、跨地域泛化能力的「通用都市微氣候代理模型（Universal Urban Microclimate Surrogate）」。

    \item \textbf{行人流量資料整合熱舒適人行動線評估：}
    結合真實世界之行人流量資料（來自手機信令、共享單車或行人計數器），對熱舒適人行動線設計目標 $f_2$ 進行人流加權修正。當前研究以均勻權重評估動線網絡各邊，未來可融入「高人流量路段優先改善熱舒適度」的需求導向最佳化。

    \item \textbf{材質與介面精細化建模：}
    現有模型以地表溫度之一階估算（式~\eqref{eq:surface_temp}）代替精確材質建模。未來可引入不同鋪面材質（瀝青、透水磚、綠化鋪面、白色反射鋪面）之熱物性資料庫，以材質節點擴充異質圖之節點類型，進一步提升模型對都市設計精細材料決策的響應能力。

    \item \textbf{現地實測驗證與感測器網路部署：}
    在新竹市典型行區部署密集 IoT UTCI 感測器網路，收集連續夏季逐時實測資料，以此進行模型物理驗證（Physical Validation）並修正現有 EnergyPlus 模擬基準與現地微氣候之系統誤差，完善研究方法論之完整性。

    \item \textbf{最佳化引擎之開發強度下限約束與三目標整合：}
    針對第~\ref{subsec:pareto_morphology}節揭露之退化解問題，於 \texttt{07\_optimization/constraints.py} 之 5 維違規向量中新增最低容積率（FAR $\ge$ FAR$_{\min}$）或建築棟數下限之可行性約束，並將人行動線熱暴露目標 $f_2$（式~\eqref{eq:walkway_exposure}）實際整合進 \texttt{FitnessEvaluator}，完成真實三目標最佳化運行以驗證多目標相互制衡能否從根本上修正單一退化方向之風險。
\end{enumerate}

\noindent\textbf{研究展望：}

本研究所建立之物理資訊時空圖神經網路框架，結合真實都市 GIS 建物擷取流程，為未來「智慧都市設計（Smart Urban Design）」工具鏈提供了一個方法論基礎。隨著感測器網路持續密集化、計算硬體持續進化，以及建築數位孿生（Building Digital Twin）技術的成熟，可以預見在五至十年內，具備物理一致性保證、能實現跨氣候跨尺度預測且錨定於真實都市紋理的多模態都市代理模型，將成為都市設計實踐的標準工具之一。本研究提出的異質圖編碼策略、物理損失約束機制、真實建物擷取方法與代理驅動多目標優化架構，可望為此演進路徑提供具體的技術參照。
